% ============================================================
%  CineTrope - Relazione di progetto per il corso di Semantic Web
%  A.A. 2025/26 - Corso di Laurea Magistrale in Informatica
% ============================================================
\documentclass[a4paper,11pt]{article}

\usepackage[utf8]{inputenc}
\usepackage[T1]{fontenc}
\usepackage[italian]{babel}
\usepackage{lmodern}
\usepackage[a4paper,top=2.4cm,bottom=2.6cm,left=2.6cm,right=2.6cm]{geometry}
\usepackage{microtype}
\usepackage{amsmath}
\usepackage{graphicx}
\usepackage{booktabs}
\usepackage{longtable}
\usepackage{array}
\usepackage{xcolor}
\usepackage{listings}
\usepackage{tikz}
\usetikzlibrary{positioning,arrows.meta}
\usepackage{caption}
\usepackage{enumitem}
\usepackage[hidelinks]{hyperref}

\hypersetup{
  pdftitle={CineTrope: un recommender semantico spiegabile},
  pdfauthor={Relazione di progetto Semantic Web},
  colorlinks=true,
  linkcolor=blue!50!black,
  urlcolor=blue!50!black,
  citecolor=blue!50!black
}

% ------------------------------------------------------------ 
% Stili per i listati
% ------------------------------------------------------------
\definecolor{ctkw}{RGB}{0,90,160}
\definecolor{ctcmt}{RGB}{0,120,60}
\definecolor{ctstr}{RGB}{150,60,20}

\lstdefinestyle{ttl}{
  basicstyle=\ttfamily\footnotesize,
  keywordstyle=\color{ctkw}\bfseries,
  commentstyle=\color{ctcmt}\itshape,
  stringstyle=\color{ctstr},
  showstringspaces=false,
  breaklines=true,
  breakatwhitespace=true,
  columns=fullflexible,
  keepspaces=true,
  frame=single,
  framesep=4pt,
  rulecolor=\color{black!25},
  morekeywords={@prefix,a,rdf:type,rdfs:Class,rdfs:subClassOf,rdfs:subPropertyOf,
    rdfs:domain,rdfs:range,rdfs:label,rdfs:comment,
    owl:Class,owl:Ontology,owl:ObjectProperty,owl:DatatypeProperty,
    owl:SymmetricProperty,owl:TransitiveProperty,owl:inverseOf,
    owl:equivalentClass,owl:equivalentProperty,owl:disjointWith,
    owl:Restriction,owl:onProperty,owl:minCardinality,owl:maxCardinality,
    owl:versionInfo,xsd:gYear,xsd:date,xsd:nonNegativeInteger,
    swrl:Imp,swrl:body,swrl:head,swrl:ClassAtom,swrl:IndividualPropertyAtom,
    swrl:classPredicate,swrl:propertyPredicate,swrl:argument1,swrl:argument2},
  alsoletter={:@},
  morecomment=[l]{\#},
}

\lstdefinestyle{sparql}{
  basicstyle=\ttfamily\footnotesize,
  keywordstyle=\color{ctkw}\bfseries,
  commentstyle=\color{ctcmt}\itshape,
  stringstyle=\color{ctstr},
  showstringspaces=false,
  breaklines=true,
  breakatwhitespace=true,
  columns=fullflexible,
  keepspaces=true,
  frame=single,
  framesep=4pt,
  rulecolor=\color{black!25},
  morekeywords={SELECT,CONSTRUCT,ASK,DESCRIBE,WHERE,PREFIX,FILTER,GROUP,BY,
    ORDER,HAVING,COUNT,DISTINCT,OPTIONAL,UNION,LIMIT,GRAPH,SERVICE,BIND,
    COALESCE,AS,DESC,ASC,STR,IN,NOT,EXISTS,MIN,MAX},
  alsoletter={:?},
  morecomment=[l]{\#},
  sensitive=false,
}

\lstdefinestyle{py}{
  language=Python,
  basicstyle=\ttfamily\footnotesize,
  keywordstyle=\color{ctkw}\bfseries,
  commentstyle=\color{ctcmt}\itshape,
  stringstyle=\color{ctstr},
  showstringspaces=false,
  breaklines=true,
  frame=single,
  framesep=4pt,
  rulecolor=\color{black!25},
  numbers=left,
  numberstyle=\tiny\color{black!40},
}

\lstdefinestyle{txt}{
  basicstyle=\ttfamily\footnotesize,
  commentstyle=\color{ctcmt}\itshape,
  breaklines=true,
  frame=single,
  framesep=4pt,
  rulecolor=\color{black!25},
  morecomment=[l]{\#},
}

\newcommand{\code}[1]{\texttt{#1}}
\newcommand{\prefix}[1]{\texttt{#1}}

\title{\textbf{CineTrope}\\[4pt]
\Large Un recommender semantico spiegabile per film,\\
basato su ontologie RDFS/OWL, named graph e ragionamento automatico\\[10pt]
\normalsize Relazione di progetto per il corso di \emph{Semantic Web}\\
Corso di Laurea Magistrale in Informatica --- A.A.\ 2025/26}
\author{}
\date{}

\begin{document}

\maketitle
\thispagestyle{empty}
\vspace{1.2cm}
\begin{center}
\begin{minipage}{0.86\textwidth}
\small
\textbf{Sommario.} 
CineTrope \`e un sistema di raccomandazione semantico per contenuti audiovisivi:
dato un film, propone cinque titoli simili e \emph{spiega} ogni suggerimento
mostrando il percorso di conoscenza che lo motiva: i tropi narrativi condivisi,
il regista, gli attori, i generi e le fonti da cui ogni informazione proviene.
Il progetto applica in modo concreto la teoria del corso di Semantic Web:
il modello RDF e le sue serializzazioni, RDF Schema per la tassonomia del
dominio, OWL 2 per disgiunzioni, propriet\`a inverse e simmetriche, equivalenze
e restrizioni, il profilo OWL RL con il reasoner \code{owlrl}, regole SWRL
eseguite da un motore forward chaining dedicato, e SPARQL per le competency
questions, per il grafo di spiegazione (\code{CONSTRUCT}) e per
l'arricchimento federato (\code{SERVICE}). La provenienza di ogni asserzione
\`e tracciata tramite named graph. La relazione approfondisce le scelte di
modellazione ontologica e illustra come ogni costrutto teorico si traduce in
una funzionalit\`a osservabile nell'applicazione.
\end{minipage}
\end{center}

\vfill
\begin{center}
\rule{0.5\textwidth}{0.4pt}
\end{center}
\newpage

\tableofcontents
\newpage

% ============================================================
\section{Introduzione}
% ============================================================

I sistemi di raccomandazione moderni sono spesso delle \emph{black-box}: un
modello, tipicamente appreso dai dati, produce un punteggio di preferenza che
l'utente pu\`o solo accettare o ignorare, senza alcuna indicazione sul
\emph{perch\'e} di un suggerimento. La visione del Semantic Web descritta a
lezione~\cite{allemang2020,antoniou2012} suggerisce un'alternativa:
rappresentare la conoscenza con un modello dichiarativo --- triple RDF
organizzate in un'ontologia --- e usare l'inferenza automatica per derivare i
risultati, in modo che ogni conclusione possa essere ricondotta ai fatti e
alle regole che l'hanno prodotta.

CineTrope nasce da questa idea. Il sistema, dato un film, propone cinque
raccomandazioni accompagnate da una \emph{spiegazione strutturata}: i tropi
narrativi condivisi (per esempio \emph{twist ending}, \emph{unreliable
narrator}, \emph{dreams}), il regista e gli attori in comune, i generi
sovrapposti e, per ogni informazione, la fonte da cui \`e stata estratta
(TMDB, DBpedia, Wikidata, dati derivati o inferenze). La spiegazione non \`e
un testo generato \emph{a posteriori}: \`e essa stessa un grafo RDF,
prodotto da una query SPARQL \code{CONSTRUCT} e visualizzato in modo
navigabile nella interfaccia web.

\subsection{Obiettivi del progetto}

Gli obiettivi, fissati in fase di progettazione a partire dal testo di
specificazione del progetto, sono i seguenti:

\begin{itemize}[nosep]
  \item costruire una \textbf{ontologia RDFS/OWL} per il dominio audiovisivo
        (opere, persone, generi, tropi, fonti), con gerarchie di classi,
        propriet\`a inverse e simmetriche, equivalenze e restrizioni;
  \item costruire un \textbf{knowledge graph} materializzato localmente a
        partire da dati pubblici (TMDB come fonte canonica, arricchimento
        biografico da DBpedia tramite query federata), separando le fonti in
        \textbf{named graph} per tracciare la provenienza di ogni tripla;
  \item usare \textbf{inferenza OWL RL} (reasoner \code{owlrl}) e
        \textbf{regole SWRL} (motore custom a forward chaining) per derivare
        conoscenza implicita: relazioni inverse e simmetriche, classificazioni
        derivate come \code{MindBender} e legami di collaborazione frequente;
  \item definire nove \textbf{competency questions} e le corrispondenti query
        SPARQL commentate, coprendo aggregazioni (\code{GROUP BY}/\code{HAVING}),
        costruzione di grafi (\code{CONSTRUCT}) e provenienza
        (\code{GRAPH ?source});
  \item fornire una \textbf{demo web} (FastAPI, Jinja2, vis.js) e una
        \textbf{CLI} che espongano catalogo, raccomandazioni, spiegazioni,
        provenienza e statistiche, senza dipendere da endpoint pubblici a
        runtime.
\end{itemize}

\subsection{Struttura della relazione}

La \textbf{Sezione~\ref{sec:dominio}} presenta il dominio di CineTrope, le
competency questions e le fonti dati. La
\textbf{Sezione~\ref{sec:ontologia}} --- cuore della relazione --- descrive
l'ontologia realizzata e le scelte di modellazione. La
\textbf{Sezione~\ref{sec:pipeline}} descrive la pipeline ETL e l'architettura
del sistema. La \textbf{Sezione~\ref{sec:reasoning}} analizza il ragionamento
automatico e i suoi risultati numerici. La \textbf{Sezione~\ref{sec:sparql}}
illustra le query SPARQL pi\`u significative, in particolare quelle di
raccomandazione e di spiegazione. La \textbf{Sezione~\ref{sec:web}} descrive
l'applicazione web e la CLI; la \textbf{Sezione~\ref{sec:valutazione}}
valuta il sistema e ne discute i limiti in modo critico.

% ============================================================
\section{CineTrope: dominio, obiettivi e dati}\label{sec:dominio}
% ============================================================

\subsection{Idea}

CineTrope \`e un \emph{recommender} semantico per contenuti audiovisivi. Il
flusso d'uso previsto dalla demo \`e:

\begin{enumerate}[nosep]
  \item l'utente sceglie un film dal catalogo;
  \item il sistema calcola cinque raccomandazioni ordinate per rilevanza;
  \item per ogni risultato mostra tropi, regista, attori e generi condivisi
        con il film di partenza, ciascuno con la propria provenienza;
  \item il grafo della spiegazione, prodotto da \code{CONSTRUCT}, \`e
        visualizzato e navigabile nell'interfaccia.
\end{enumerate}

La differenza rispetto a un filtro collaborativo \`e paradigmatica: la
similarit\`a non \`e una distanza in uno spazio latente ma un
\emph{percorso nel grafo di conoscenza}, che l'utente pu\`o ispezionare.
Il punteggio \`e un'euristica trasparente, non un modello appreso.

\subsection{Competency questions}

Seguendo la metodologia del corso, prima della modellazione sono state
definite le domande a cui il sistema deve saper rispondere. Le nove
\emph{competency questions} (CQ) di CineTrope, implementate in
\code{cinetrope/queries.py} ed eseguibili dalla UI, dall'API e dalla CLI,
sono riportate in tabella~\ref{tab:cq}.

\begin{table}[ht]
\centering
\caption{Le competency questions di CineTrope.}
\label{tab:cq}
\small
\begin{tabular}{@{}clp{8.6cm}@{}}
\toprule
\textbf{ID} & \textbf{Titolo} & \textbf{Domanda} \\
\midrule
cq1 & Tropi condivisi & Quali film condividono almeno due tropi con
\emph{Inception}? \\
cq2 & Regista e tropo & Quali film sono diretti da persone che hanno
lavorato su opere con lo stesso tropo? \\
cq3 & MindBender & Quali film sono classificati come \code{MindBender}
dalle regole SWRL? \\
cq4 & Provenienza & Da quale fonte proviene l'asserzione ``Inception usa
\code{dreams}''? \\
cq5 & Raccomandazioni & Quali film condividono tropi, regista, attori o
generi con un film scelto? \\
cq6 & Collaborazioni & Quali persone sono collegate da
\code{hasFrequentCollaborator}? \\
cq7 & Gerarchia generi & Quali generi sono sottogeneri di altri generi? \\
cq8 & Fonti & Quante triple asserisce ciascun named graph? \\
cq9 & Co-star & Quali attori hanno recitato insieme? \\
\bottomrule
\end{tabular}
\end{table}

Le CQ guidano direttamente le scelte ontologiche: la cq3 richiede una classe
derivata (\`e una regola, non un dato); la cq4 richiede la provenienza per
asserzione (named graph); la cq9 richiede una propriet\`a simmetrica; la cq6
richiede il ponte fra aggregazione ETL e regola SWRL.

\subsection{Fonti dati e materializzazione}

Le fonti esterne del progetto sono:

\begin{itemize}[nosep]
  \item \textbf{TMDB}: fonte canonica del catalogo. Lo script
        \code{scripts/import\_tmdb.py} importa 500 film con titolo, anno,
        regia, primi 10 attori accreditati, generi e \emph{keyword}; ogni
        keyword diventa un individuo della classe \code{Tropo} e gli ID
        vengono stabilizzati come \code{TMDB<id>} per evitare duplicati fra
        importazioni;
  \item \textbf{DBpedia}: arricchimento federato delle date di nascita dei
        registi tramite la query \code{SERVICE} descritta nella
        Sezione~\ref{sec:sparql}; il risultato viene
        \emph{materializzato} in \code{data/raw/} (file
        \code{dbpedia\_birthdates.json}) e ricaricato nel named graph
        dedicato;
  \item \textbf{Wikidata}: metadati biografici inclusi nel dataset locale
        \code{data/raw/people.json};
  \item \textbf{analisi locale}: gerarchia dei generi curata manualmente in
        \code{data/raw/genres.json}.
\end{itemize}

La decisione architetturale fondamentale \`e la \textbf{materializzazione
locale}: nessun endpoint pubblico viene contattato a runtime. Le query
federate sono usate \emph{solo} durante l'ETL, come raccomandato dal testo
di progetto, per rendere la demo rapida, ripetibile e indipendente
dall'indisponibilit\`a dei servizi pubblici. Inoltre, come da indicazioni
del corso, si \`e preferito un dataset \emph{piccolo ma coerente} a uno
ampio ma rumoroso: la qualit\`a delle raccomandazioni dipende dalla
coerenza del grafo, non dalla sua dimensione.

% ============================================================
\section{L'ontologia CineTrope}\label{sec:ontologia}
% ============================================================

L'ontologia \`e serializzata in Turtle in
\code{data/ontology/cinetrope.ttl} e dichiarata come
\code{owl:Ontology} nel namespace \code{http://example.org/cinetrope\#}
(prefisso \code{:}). Questa sezione ne descrive i componenti e le scelte di
modellazione, mettendole in relazione con la teoria.

\subsection{Classi}

La tassonomia delle classi \`e riassunta in
figura~\ref{fig:classLists} e nel listato~\ref{lst:classi}.

\begin{figure}[ht]
\centering
\begin{tikzpicture}[
  every node/.style={draw, rounded corners=2pt, align=center, font=\small,
    inner sep=4pt, minimum height=1.5em},
  sub/.style={-{Latex[length=2mm]}, thick},
  label/.style={font=\scriptsize\itshape, midway, right=2pt}
]
\node (oa) {\textbf{OperaAudiovisiva}};
\node[below left=1.05cm and 1.5cm of oa] (film) {Film};
\node[below right=1.05cm and 1.5cm of oa] (tv) {SerieTV};
\node[below=1.05cm of film] (mb) {MindBender};
\node[below right=1.05cm and 0.4cm of film] (persona) {\textbf{Persona}};
\node[below left=1.05cm and 0.4cm of persona] (regista) {Regista};
\node[below right=1.05cm and 0.4cm of persona] (attore) {Attore};
\draw[sub] (film) -- node[label] {subClassOf} (oa);
\draw[sub] (tv) -- node[label, left=2pt] {subClassOf} (oa);
\draw[sub] (mb) -- node[label] {subClassOf} (film);
\draw[sub] (regista) -- node[label] {subClassOf} (persona);
\draw[sub] (attore) -- node[label] {subClassOf} (persona);
\end{tikzpicture}
\caption{Tassonomia principale dell'ontologia CineTrope.}
\label{fig:classLists}
\end{figure}

\begin{lstlisting}[style=ttl, caption={Frammento di
\code{data/ontology/cinetrope.ttl}: classi principali.}, label={lst:classi}]
@prefix : <http://example.org/cinetrope#> .
@prefix owl: <http://www.w3.org/2002/07/owl#> .
@prefix rdfs: <http://www.w3.org/2000/01/rdf-schema#> .

:OperaAudiovisiva a owl:Class ;
    rdfs:label "Opera audiovisiva"@it .

:Film a owl:Class ;
    rdfs:subClassOf :OperaAudiovisiva ;
    owl:disjointWith :SerieTV ;
    rdfs:label "Film"@it .

:SerieTV a owl:Class ;
    rdfs:subClassOf :OperaAudiovisiva .

:Movie a owl:Class ;
    owl:equivalentClass :Film .

:Regista a owl:Class ; rdfs:subClassOf :Persona .
:Attore a owl:Class ; rdfs:subClassOf :Persona .

:MindBender a owl:Class ;
    rdfs:subClassOf :Film ;
    rdfs:comment "Classe derivata dalla regola SWRL: un film con
      UnreliableNarrator e TwistEnding e' un MindBender."@it .
\end{lstlisting}

Le scelte significative sono:

\begin{itemize}
  \item \textbf{Superclasse \code{OperaAudiovisiva}}: il dominio include
        film e serie TV; le propriet\`a comuni (\code{hasGenre},
        \code{directedBy}, \code{starring}, \code{hasTrope},
        \code{releaseYear}) hanno domain su di essa, cos\`icch\'e la stessa
        relazione \`e riusabile per entrambi i tipi di opera senza
        duplicare vocaboli (\emph{union pattern} su classi concrete);
  \item \textbf{\code{Film owl:disjointWith SerieTV}}: l'assioma di
        disgiunzione serve esattamente allo scopo visto a lezione ---
        \emph{rilevare} classificazioni incompatibili nei dati importati;
        il controllo \code{film\_and\_series} della validazione
        (Sezione~\ref{sec:validazione}) verifica che nessuna risorsa sia
        entrambe le cose;
  \item \textbf{allineamento terminologico}: \code{Movie} \`e dichiarata
        \code{owl:equivalentClass} di \code{Film}, e \code{TVSeries} di
        \code{SerieTV}: \`e l'applicazione diretta della
        conciliazione di termini; un vocabolario esterno che parla di
        \code{Movie} viene ricondotto automaticamente a \code{Film};
  \item \textbf{classi di ruolo}: \code{Regista} e \code{Attore} sono
        sottoclassi di \code{Persona}; un individuo pu\`o avere entrambi i
        ruoli (non sono disgiunte), cosa frequente nei dati reali;
  \item \textbf{\code{MindBender} \`e una classe \emph{derivata}}: non \`e
        mai asserita nei dati, ma \`e conclusa dalla regola SWRL
        corrispondente (Sezione~\ref{sec:swrl}). Essere sottoclasse di
        \code{Film} significa che ogni istanza inferita gode anche di tutte
        le propriet\`a dei film.
\end{itemize}

Completano il vocabolario di classi: \code{Genere}, \code{Tropo},
\code{ElementoDescrittivo}, \code{Scena}, \code{Fonte}. La classe
\code{ElementoDescrittivo} modella le keyword TMDB non narrative
(ambientazioni, oggetti, contesti): nell'importazione corrente tutte le
keyword sono trattate come \code{Tropo} per massimizzare la copertura delle
raccomandazioni, ma la distinzione resta disponibile nell'ontologia per una
curation futura --- esempio di come l'ontologia codifichi anche la
\emph{visione} del dominio, non solo i dati presenti.

\subsection{Propriet\`a}

Il listato~\ref{lst:props} mostra la dichiarazione tipica delle
propriet\`a; la tabella~\ref{tab:props} riporta per ciascuna domain, range
e caratteristiche.

\begin{lstlisting}[style=ttl, caption={Frammento dell'ontologia:
propriet\`a inverse, simmetriche ed equivalente.}, label={lst:props}]
:directedBy a owl:ObjectProperty ;
    owl:inverseOf :directorOf ;
    rdfs:domain :OperaAudiovisiva ;
    rdfs:range :Regista .

:coStarsWith a owl:ObjectProperty , owl:SymmetricProperty ;
    rdfs:domain :Attore ;
    rdfs:range :Attore .

:usesTrope a owl:ObjectProperty ;
    owl:equivalentProperty :hasTrope ;
    rdfs:domain :OperaAudiovisiva ;
    rdfs:range :Tropo .

:releaseYear a owl:DatatypeProperty ;
    rdfs:domain :OperaAudiovisiva ;
    rdfs:range xsd:gYear .
\end{lstlisting}

\begin{table}[ht]
\centering
\caption{Propriet\`a dell'ontologia CineTrope.}
\label{tab:props}
\footnotesize
\begin{tabular}{@{}llll@{}}
\toprule
\textbf{Propriet\`a} & \textbf{Domain} & \textbf{Range} & \textbf{Caratteristiche} \\
\midrule
\code{hasGenre} & OperaAudiovisiva & Genere & --- \\
\code{subGenreOf} & Genere & Genere & transitiva \\
\code{directedBy} & OperaAudiovisiva & Regista & inversa di \code{directorOf} \\
\code{directorOf} & Regista & OperaAudiovisiva & inversa di \code{directedBy} \\
\code{starring} & OperaAudiovisiva & Attore & inversa di \code{actedIn} \\
\code{actedIn} & Attore & OperaAudiovisiva & inversa di \code{starring} \\
\code{coStarsWith} & Attore & Attore & simmetrica \\
\code{hasTrope} & OperaAudiovisiva & Tropo & inversa di \code{tropeUsedIn} \\
\code{usesTrope} & OperaAudiovisiva & Tropo & equivalente a \code{hasTrope} \\
\code{tropeUsedIn} & Tropo & OperaAudiovisiva & inversa di \code{hasTrope} \\
\code{hasElement} & OperaAudiovisiva & ElementoDescrittivo & inv.\ \code{elementUsedIn} \\
\code{hasScene} & OperaAudiovisiva & Scena & inversa di \code{sceneIn} \\
\code{annotatedWith} & Scena & Tropo & --- \\
\code{statedBy} & Scena & Fonte & --- \\
\code{collaboratedWith} & Persona & Persona & simmetrica \\
\code{frequentCollaboration} & Persona & Persona & simmetrica, da ETL \\
\code{hasFrequentCollaborator} & Persona & Persona & simmetrica, inferita SWRL \\
\code{releaseYear} & OperaAudiovisiva & \code{xsd:gYear} & datatype property \\
\code{birthDate} & Persona & \code{xsd:date} & datatype property \\
\midrule
\multicolumn{4}{@{}l@{}}{\emph{vocabolario del grafo di spiegazione}} \\
\code{recommends} & Film & Film & --- \\
\code{recommendedBecause} & Film & Tropo & --- \\
\code{sharesTropeWith} & Film & Film & --- \\
\code{sharesDirectorWith} & Film & Film & --- \\
\code{sharesActorWith} & Film & Film & --- \\
\code{sharesGenreWith} & Film & Film & --- \\
\code{assertedIn} & rdfs:Resource & Fonte & --- \\
\bottomrule
\end{tabular}
\end{table}

Ogni costrutto ha una funzione precisa nell'applicazione:

\begin{itemize}
  \item \textbf{coppie inverse} (\code{directedBy}/\code{directorOf},
        \code{starring}/\code{actedIn}, \code{hasTrope}/\code{tropeUsedIn},
        \code{hasElement}/\code{elementUsedIn},
        \code{hasScene}/\code{sceneIn}): la lezione sul ``matrimonio di
        Shakespeare'' ha mostrato il problema: senza inverso dichiarato, la
        query \code{SELECT ?film WHERE \{?film :directorOf :Nolan\}} non
        trova nulla. CineTrope dichiara gli inversi in modo che il reasoner
        \emph{completi} il grafo in entrambe le direzioni; i contatori della
        Sezione~\ref{sec:reasoning} mostrano 554 \code{directorOf} e 4.875
        \code{actedIn} \emph{inferite} a partire dai dati direzionali
        importati;
  \item \textbf{propriet\`a simmetriche} (\code{coStarsWith},
        \code{collaboratedWith}, \code{frequentCollaboration},
        \code{hasFrequentCollaborator}): il fatto ``due attori recitano
        insieme'' \`e per natura simmetrico; dichiarare
        \code{owl:SymmetricProperty} evita di asserire entrambe le direzioni
        e --- come dettagliato nella Sezione~\ref{sec:reasoning} --- il
        progetto lascia \code{coStarsWith} monodirezionale nei dati derivati
        proprio per rendere l'inferenza OWL visibile e verificabile;
  \item \textbf{propriet\`a equivalente}: \code{usesTrope
        owl:equivalentProperty hasTrope} \`e l'applicazione della
        conciliazione di termini vista a lezione (il caso
        \code{borrows}/\code{checkedOut} delle biblioteche): un vocabolario
        esterno che usi \code{usesTrope} viene fuso con quello interno senza
        duplicare i dati;
  \item \textbf{\code{subGenreOf} transitiva}: la gerarchia dei generi
        (\emph{ScienceFiction subGenreOf Fantascienza}, ecc.) beneficia della
        transitivit\`a: le query che attraversano la gerarchia non devono
        enumerare i livelli intermedi.
\end{itemize}

\subsection{I generi come individui}

Una scelta di modellazione discussa esplicitamente: i generi sono
\textbf{individui} della classe \code{Genere}, collegati da
\code{subGenreOf}, e non classi OWL. L'alternativa --- il \emph{punning}
OWL 2, in cui un genere come \code{ScienceFiction} sarebbe
contemporaneamente una \code{owl:Class} (per classificare i film) e un
individuo di \code{Genere} (per avere etichette, gerarchia e metadati) ---
\`e stata valutata e scartata, perch\'e la modellazione a individui basta
agli scopi del progetto: la relazione film--genere \`e una normale propriet\`a
\code{hasGenre}, la gerarchia \`e una propriet\`a transitiva tra individui, e
le query restano uniformi. Durante l'ETL l'identificatore TMDB
\code{SciFi} viene normalizzato in \code{ScienceFiction}, mentre
l'allineamento tra terminologie \`e delegato a
\code{owl:equivalentClass}/\code{owl:equivalentProperty} dove serve. Questa
\`e un'applicazione consapevole del principio ``la conciliazione dei
termini si fa con assiomi, non rinominando i dati''.

\subsection{Restrizioni OWL e validazione}\label{sec:validazione}

L'ontologia dichiara due restrizioni di cardinalit\`a:

\begin{lstlisting}[style=ttl, caption={Restrizioni di cardinalit\`a su
\code{Film}.}]
:Film rdfs:subClassOf [
    a owl:Restriction ;
    owl:onProperty :directedBy ;
    owl:minCardinality "1"^^xsd:nonNegativeInteger
] .

:Film rdfs:subClassOf [
    a owl:Restriction ;
    owl:onProperty :releaseYear ;
    owl:maxCardinality "1"^^xsd:nonNegativeInteger
] .
\end{lstlisting}

Come discusso a lezione, il reasoner serve a \emph{individuare} le
incoerenze, non a correggerle. C'\`e per\`o un limite tecnico preciso,
dichiarato onestamente nel progetto: il profilo \textbf{OWL RL} non copre
tutte le cardinalit\`a di OWL 2 DL (in particolare non le cardinalit\`a
qualificate su classi arbitrarie), quindi il progetto affianca alle
restrizioni dei \textbf{controlli SPARQL espliciti} eseguiti
all'ETL e salvati nel report di build:

\begin{itemize}[nosep]
  \item \code{missing\_director}: film senza alcun \code{directedBy}
        (violazione del vincolo min);
  \item \code{missing\_year}: film senza \code{releaseYear};
  \item \code{multiple\_years}: film con pi\`u anni di uscita (violazione
        del vincolo max, individuata con \code{GROUP BY}/\code{HAVING});
  \item \code{film\_and\_series}: risorse sia \code{Film} sia
        \code{SerieTV} (violazione della disgiunzione);
  \item \code{owl\_nothing}: individui classificati \code{owl:Nothing} dal
        reasoner, cio\`e rese inconsistenti da qualche assioma.
\end{itemize}

Nel dataset corrente tutti i controlli restituiscono zero violazioni. Si
noti la coerenza metodologica con l'OWA: i controlli \emph{segnalano} dati
mancanti o ambigui senza sostituirli con valori inventati.

\subsection{Named graph e provenienza}

Le fonti sono modellate come individui della classe \code{Fonte} e usate
come nomi dei relativi named graph. La tabella~\ref{tab:graphs} riporta
l'organizzazione del dataset (un file TriG con quadruple).

\begin{table}[ht]
\centering
\caption{Named graph del dataset CineTrope.}
\label{tab:graphs}
\small
\begin{tabular}{@{}lp{9.8cm}@{}}
\toprule
\textbf{Named graph} & \textbf{Contenuto} \\
\midrule
\code{:Ontology} & schema RDFS/OWL e vocabolario (cinetrope.ttl) \\
\code{:TMDB} & titoli, anni, regia, cast, generi e keyword importati \\
\code{:Wikidata} & metadati strutturati delle persone (date di nascita) \\
\code{:DBpedia} & date di nascita arricchite via \code{SERVICE} \\
\code{:TvTropes} & riservato alle annotazioni narrative esterne \\
\code{:LocalAnalysis} & generi e gerarchie curati nel dataset locale \\
\code{:Derived} & co-star e collaborazioni calcolate con \code{CONSTRUCT} \\
\code{:Inferred} & chiusura OWL RL e conseguenze SWRL \\
\bottomrule
\end{tabular}
\end{table}

La provenienza \`e interrogabile direttamente, senza reificazione, come
illustrato a lezione:

\begin{lstlisting}[style=sparql, caption={Provenienza di un'asserizione
tramite named graph (\code{05\_provenance\_assertion.rq}).}]
PREFIX : <http://example.org/cinetrope#>
PREFIX rdfs: <http://www.w3.org/2000/01/rdf-schema#>

SELECT ?source ?sourceLabel WHERE {
  GRAPH ?source {
    :TMDB27205 :hasTrope :dreams .
  }
  OPTIONAL { ?source rdfs:label ?sourceLabel }
}
\end{lstlisting}

Il grafo di spiegazione collega poi i fatti alle fonti con la propriet\`a
\code{assertedIn}, cos\`icch\'e anche la visualizzazione vis.js espone la
provenienza di ogni arco. Il confronto fra sorgenti --- tema della lezione
sull'integrazione --- \`e qui ridotto a una sola clausola
\code{GRAPH ?source}.

% ============================================================
\section{Pipeline dati e architettura}\label{sec:pipeline}
% ============================================================

L'architettura segue la pipeline consigliata dal testo di progetto
(API/ETL $\rightarrow$ triple RDF $\rightarrow$ ontologia $\rightarrow$
reasoner $\rightarrow$ store $\rightarrow$ interfaccia). La
figura~\ref{fig:arch} mostra il flusso di costruzione e la
tabella~\ref{tab:moduli} i componenti software.

\begin{figure}[ht]
\centering
\begin{tikzpicture}[
  node distance=0.55cm and 1.4cm,
  box/.style={draw, rounded corners=2pt, align=center, font=\scriptsize,
    inner sep=5pt, minimum width=2.6cm, minimum height=1.6em},
  data/.style={box, fill=black!4},
  proc/.style={box},
  store/.style={box, fill=black!10},
  app/.style={box, fill=black!4},
  arr/.style={-{Latex[length=2mm]}, thick}
]
\node[data] (raw) {data/raw/*.json\\TMDB, DBpedia,\\Wikidata, generi};
\node[proc, right=of raw] (etl) {ETL rdflib\\(cinetrope/etl.py)};
\node[data, right=of etl] (ont) {cinetrope.ttl\\+ regole SWRL};
\node[store, below=of etl] (ds) {Dataset rdflib\\named graph per fonte};
\node[proc, right=of ds] (derive) {CONSTRUCT derivati\\co-star, collaborazioni};
\node[proc, below=of ds] (owlrl) {OWL RL (owlrl)\\+ SWRL custom};
\node[store, below=of owlrl] (trig) {store.trig\\+ build\_report.json};
\node[app, below=of derive] (api) {FastAPI + Jinja2\\API JSON};
\node[app, below=of api] (ui) {vis.js / CLI\\UI web};

\draw[arr] (raw) -- (etl);
\draw[arr] (ont) -- (ds);
\draw[arr] (etl) -- (ds);
\draw[arr] (ds) -- (derive);
\draw[arr] (derive) -- (owlrl);
\draw[arr] (owlrl) -- (trig);
\draw[arr] (trig) -- (api);
\draw[arr] (api) -- (ui);
\end{tikzpicture}
\caption{Pipeline di costruzione del knowledge graph e stack applicativo.}
\label{fig:arch}
\end{figure}

\begin{table}[ht]
\centering
\caption{Componenti software del progetto.}
\label{tab:moduli}
\small
\begin{tabular}{@{}lp{10.4cm}@{}}
\toprule
\textbf{File} & \textbf{Ruolo} \\
\midrule
\code{cinetrope/etl.py} & costruisce il dataset RDF e i named graph \\
\code{cinetrope/reasoning.py} & OWL RL con \code{owlrl}, inferenze e validazione \\
\code{cinetrope/swrl.py} & parser ed esecuzione delle regole SWRL \\
\code{cinetrope/queries.py} & caricamento, rendering e competency questions \\
\code{cinetrope/store.py} & API applicativa: catalogo, raccomandazioni, spiegazioni, provenienza \\
\code{cinetrope/web.py} & rotte FastAPI e template \\
\code{cinetrope/cli.py} & interfaccia da terminale \\
\code{cinetrope/tmdb.py} & client TMDB e importazione keyword/cast \\
\code{scripts/import\_tmdb.py} & importazione e aggiornamento del catalogo TMDB \\
\code{scripts/enrich\_birthdates.py} & arricchimento federato da DBpedia \\
\code{scripts/build\_store.py} & build completo: ETL + reasoner + validazione \\
\bottomrule
\end{tabular}
\end{table}

\subsection{ETL: dai JSON alle triple}

La funzione \code{build\_catalog} in \code{cinetrope/etl.py} legge i dati
grezzi e popola un \code{Dataset} rdflib in modalit\`a \code{default\_union}:

\begin{itemize}[nosep]
  \item \code{\_load\_ontology} carica \code{cinetrope.ttl} nel grafo
        \code{:Ontology};
  \item \code{\_add\_source\_facts} scrive generi e gerarchia in
        \code{:LocalAnalysis} e le keyword TMDB in \code{:TMDB} (ogni
        keyword diventa un individuo \code{:Tropo} con \code{rdfs:label}
        in italiano);
  \item \code{\_add\_people} crea le persone con ID locali derivati dal nome
        (\code{slugify}) e i tipi \code{:Persona}, \code{:Regista},
        \code{:Attore}; le date di nascita finiscono nel grafo della fonte
        (\code{:DBpedia} se arricchite, \code{:Wikidata} altrimenti);
  \item \code{\_add\_films} crea un individuo per opera con ID stabile
        \code{TMDB<id>}, tipi \code{:Film}/\code{:SerieTV},
        \code{rdfs:label}, \code{:releaseYear} tipato \code{xsd:gYear},
        e le relazioni \code{:hasGenre}, \code{:directedBy},
        \code{:starring}, \code{:hasTrope};
  \item \code{\_add\_derived\_relations} calcola con \code{CONSTRUCT} le
        co-star e le collaborazioni e, con una query aggregata, individua
        le collaborazioni frequenti.
\end{itemize}

Il punto metodologicamente interessante \`e la divisione del lavoro fra
\textbf{ETL, SPARQL e regole}: le relazioni che dipendono da conteggi
(\`e il limite delle aggregazioni in SWRL) sono calcolate da query
\code{CONSTRUCT} con \code{GROUP BY}/\code{HAVING} e materializzate nel
grafo \code{:Derived}; le conseguenze logiche sono demandate al reasoner.
Un esempio \`e \code{derive\_co\_stars.rq}:

\begin{lstlisting}[style=sparql, caption={Materializzazione ETL delle
co-star (una sola direzione, per mostrare l'inferenza OWL).}]
PREFIX : <http://example.org/cinetrope#>

CONSTRUCT { ?actor1 :coStarsWith ?actor2 . }
WHERE {
  ?work :starring ?actor1, ?actor2 .
  FILTER (STR(?actor1) < STR(?actor2))
}
\end{lstlisting}

La condizione \code{STR(?a) < STR(?b)} garantisce una sola tripla per
coppia; la direzione inversa --- come vedremo --- sar\`a inferita dalla
simmetria dichiarata in OWL. Le collaborazioni frequenti usano invece
l'aggregazione:

\begin{lstlisting}[style=sparql, caption={Rilevamento delle collaborazioni
frequenti: il conteggio che SWRL non pu\`o fare.}, label={lst:frequent}]
PREFIX : <http://example.org/cinetrope#>

CONSTRUCT { ?p1 :frequentCollaboration ?p2 . }
WHERE {
  SELECT ?p1 ?p2 (COUNT(DISTINCT ?work) AS ?works)
  WHERE {
    ?p1 :collaboratedWith ?p2 .
    ?work :directedBy|:starring ?p1 , ?p2 .
  }
  GROUP BY ?p1 ?p2
  HAVING (COUNT(DISTINCT ?work) > 3)
}
\end{lstlisting}

La regola SWRL completa poi il quadro collegando gli individui con
\code{hasFrequentCollaborator} (Sezione~\ref{sec:swrl}).

% ============================================================
\section{Ragionamento automatico}\label{sec:reasoning}
% ============================================================

\subsection{Chiusura OWL RL}

Il ragionamento \`e implementato in \code{cinetrope/reasoning.py} con la
libreria \code{owlrl}:

\begin{lstlisting}[style=py, caption={Esecuzione della chiusura OWL RL e
delle regole SWRL (cinetrope/reasoning.py).}]
def run_reasoning(dataset):
    union = _union_graph(dataset)      # unione di tutti i grafi espliciti
    explicit = set(union)

    _run_owlrl(union)                  # prima chiusura OWL RL + RDFS
    swrl_stats = apply_swrl(union, load_rule_graph())
    _run_owlrl(union)                  # seconda chiusura dopo le regole

    inferred = {                        # solo triple che toccano il ns CineTrope
        (s, p, o) for s, p, o in (set(union) - explicit)
        if _is_cinetrope(s) or _is_cinetrope(p) or _is_cinetrope(o)
        and not (p == OWL.sameAs and s == o)
    }
    dataset.graph(CT.Inferred).add(...)  # salvate nel grafo :Inferred
\end{lstlisting}

La strategia \`e in tre passi:

\begin{enumerate}[nosep]
  \item i named graph espliciti vengono fusi in un grafo di lavoro
        (\emph{merging}, favorito dalla natura stessa di RDF);
  \item \code{DeductiveClosure(OWLRL\_Semantics)} con
        \code{rdfs\_closure=True} espande la chiusura OWL RL + RDFS ---
        ovvero \emph{materializza} le conseguenze delle propriet\`a
        inverse, simmetriche, equivalenti, delle gerarchie di classi e
        propriet\`a e dei domain/range;
  \item le regole SWRL vengono applicate sul grafo gi\`a espanso, e una
        seconda passata di OWL RL stabilizza le eventuali conseguenze
        combinate.
\end{enumerate}

Le nuove triple, filtrate a quelle che coinvolgono il namespace CineTrope
(ed escludendo le tautologie \code{owl:sameAs} introdotte dalla chiusura),
vengono scritte nel named graph \code{:Inferred}. La separazione
\emph{dati espliciti / dati inferiti} in named graph diversi \`e essa
stessa un'applicazione del pattern di provenienza: si pu\`o sempre
distinguere ci\`o che \`e stato asserito da ci\`o che \`e stato dedotto.

\subsection{Risultati dell'inferenza}

Sul dataset corrente (500 film, 3.929 persone, 3.514 tropi, 21 generi,
161.843 triple totali) la chiusura produce \textbf{52.608 triple
inferite}, ripartite come in tabella~\ref{tab:inference}.

\begin{table}[ht]
\centering
\caption{Ripartizione delle 52.608 triple inferite per predicato
(build\_report.json).}
\label{tab:inference}
\small
\begin{tabular}{@{}lr@{\hspace{1.4cm}}lr@{}}
\toprule
\textbf{Predicato} & \textbf{Triple} & \textbf{Predicato} & \textbf{Triple} \\
\midrule
\code{coStarsWith} & 20.807 & \code{directorOf} & 554 \\
\code{rdf:type} & 10.519 & \code{subClassOf} & 110 \\
\code{tropeUsedIn} & 7.727 & \code{hasFrequentCollaborator} & 92 \\
\code{usesTrope} & 7.727 & \code{domain}/\code{range} & 54 / 49 \\
\code{actedIn} & 4.875 & \code{equivalentClass} & 35 \\
 & & \code{subPropertyOf} & 30 \\
\bottomrule
\end{tabular}
\end{table}

La lettura della tabella \`e una lettura della teoria in azione:

\begin{itemize}
  \item \textbf{20.807 \code{coStarsWith}} inferite: l'ETL materializza una
        sola direzione per ogni coppia di co-star; la simmetria dichiarata
        in OWL produce l'altra direzione. \`E la dimostrazione pratica del
        costrutto \code{owl:SymmetricProperty} visto a lezione;
  \item \textbf{7.727 \code{tropeUsedIn} e 7.727 \code{usesTrope}}: la prima
        \`e l'inverso di \code{hasTrope}, la seconda la conseguenza
        dell'\code{owl:equivalentProperty}. Le 7.726 asserzioni
        \code{hasTrope} importate da TMDB generano, senza alcuna duplicazione
        manuale, il doppio del valore interrogabile;
  \item \textbf{4.875 \code{actedIn} e 554 \code{directorOf}}: gli inversi
        di \code{starring} e \code{directedBy};
  \item \textbf{10.519 \code{rdf:type}}: propagazione lungo le gerarchie
        (\code{Regista}/\code{Attore} $\rightarrow$ \code{Persona},
        \code{Film} $\rightarrow$ \code{OperaAudiovisiva}) e classificazione
        automatica tramite \code{rdfs:domain}/\code{rdfs:range};
  \item \textbf{92 \code{hasFrequentCollaborator}}: conseguenze della regola
        SWRL, non di assiomi OWL.
\end{itemize}

\subsection{Il motore SWRL custom}\label{sec:swrl}

Le regole SWRL sono dichiarate in
\code{data/ontology/cinetrope-swrl.n3}, in serializzazione RDF del
metamodello SWRL (\code{swrl:Imp}, \code{swrl:body}, \code{swrl:head}, con
atomi \code{swrl:ClassAtom} e \code{swrl:IndividualPropertyAtom}). Il
formato N3 \`e stato scelto perch\'e il parser Turtle di rdflib non
supporta le variabili \code{?x} come termini RDF. La regola principale:

\begin{lstlisting}[style=ttl, caption={La regola MindBenderRule in
\code{cinetrope-swrl.n3} (metamodello SWRL in RDF).}]
@prefix : <http://example.org/cinetrope#> .
@prefix swrl: <http://www.w3.org/2003/11/swrl#> .

:MindBenderRule a swrl:Imp ;
    rdfs:label "Un film con narratore inaffidabile e finale a sorpresa
                e' un MindBender"@it ;
    swrl:body (
        [ a swrl:ClassAtom ;
          swrl:classPredicate :Film ;
          swrl:argument1 ?f ]
        [ a swrl:IndividualPropertyAtom ;
          swrl:propertyPredicate :hasTrope ;
          swrl:argument1 ?f ;
          swrl:argument2 :UnreliableNarrator ]
        [ a swrl:IndividualPropertyAtom ;
          swrl:propertyPredicate :hasTrope ;
          swrl:argument1 ?f ;
          swrl:argument2 :TwistEnding ]
    ) ;
    swrl:head (
        [ a swrl:ClassAtom ;
          swrl:classPredicate :MindBender ;
          swrl:argument1 ?f ]
    ) .
\end{lstlisting}

Il file parallelo \code{cinetrope-rules.swrl} riporta la stessa regola in
sintassi leggibile, corrispondente a quella vista a lezione:

\begin{lstlisting}[style=txt, caption={La stessa regola in sintassi
SWRL leggibile.}]
Film(?f) ^ hasTrope(?f, :UnreliableNarrator)
       ^ hasTrope(?f, :TwistEnding)
  -> MindBender(?f)

Persona(?p1) ^ Persona(?p2)
       ^ collaboratedWith(?p1, ?p2)
       ^ frequentCollaboration(?p1, ?p2)
  -> hasFrequentCollaborator(?p1, ?p2)
\end{lstlisting}

Il motore (\code{cinetrope/swrl.py}) implementa un \textbf{forward
chaining} fino al punto fisso (massimo 20 iterazioni): per ogni regola,
gli atomi del corpo vengono \emph{abbinati} al grafo con un algoritmo di
unificazione (\code{\_match\_atom}) che propaga i binding delle variabili
da un atomo al successivo; ogni binding completo istanzia la testa e
aggiunge le triple mancanti. Il corpo della seconda regola mostra la
collaborazione fra teoria e pratica del corso: \code{collaboratedWith} \`e
materiale derivato dall'ETL, \code{frequentCollaboration} \`e asserita
dall'aggregazione SPARQL del listato~\ref{lst:frequent}, e la regola
\emph{collega} i due mondi producendo \code{hasFrequentCollaborator}.
L'esecuzione corrente produce 92 inferenze.

Un limite dichiarato del motore: il supporto si ferma agli atomi di classe
(\code{ClassAtom}) e di propriet\`a (\code{IndividualPropertyAtom}) --- non
\`e un implementazione completa di SWRL (niente built-in, niente atomi
\code{SameAs}/\code{DifferentFrom}). La scelta \`e coerente con l'uso
didattico: poche regole, semanticamente trasparenti, verificate dai test.

% ============================================================
\section{SPARQL al lavoro: raccomandazioni e spiegazioni}\label{sec:sparql}
% ============================================================

Le tredici query del progetto sono in \code{data/queries/}, commentate.
Questa sezione analizza le due pi\`u importanti e i pattern ricorrenti.

\subsection{La query di raccomandazione}

La query \code{03\_recommendations.rq} calcola il punteggio combinato. Il
problema ingegneristico --- evitare prodotti cartesiani fra le quattro
dimensioni --- \`e risolto con quattro sottoquery \code{OPTIONAL} aggregate
e unite sul film:

\begin{lstlisting}[style=sparql, caption={Struttura della query di
raccomandazione (semplificata).}]
PREFIX : <http://example.org/cinetrope#>

SELECT ?film ?title ?year
       (COALESCE(?sharedTropes, 0) AS ?sharedTropes)
       ... (COALESCE ... per registi, attori, generi) ...
       (COALESCE(?sharedTropes,0) * {{TROPE_WEIGHT}}
        + COALESCE(?sharedDirectors,0) * {{DIRECTOR_WEIGHT}}
        + COALESCE(?sharedActors,0) * {{ACTOR_WEIGHT}}
        + COALESCE(?sharedGenres,0) * {{GENRE_WEIGHT}} AS ?score)
WHERE {
  ?film a :Film ; rdfs:label ?title .
  FILTER (?film != <{{SEED}}>)
  OPTIONAL { ?film :releaseYear ?year }

  OPTIONAL {
    SELECT ?film (COUNT(DISTINCT ?trope) AS ?sharedTropes) WHERE {
      <{{SEED}}> :hasTrope ?trope .
      ?film :hasTrope ?trope .
      FILTER (?film != <{{SEED}}>)
    } GROUP BY ?film
  }
  ... tre OPTIONAL analoghi per directedBy, starring, hasGenre ...
}
ORDER BY DESC(?score) DESC(?sharedTropes) ?title
LIMIT {{LIMIT}}
\end{lstlisting}

Ogni sottoquery conta le entit\`a condivise con il film seme; la
\code{COALESCE} gestisce i film senza sovrapposizioni (il valore
\code{UNBOUND} \`e il comportamento previsto di \code{OPTIONAL}, non un
errore). I \emph{pesi} sono parametri del template, passati
dall'applicazione in base al preset scelto dall'utente
(tabella~\ref{tab:weights}). Grazie al layer applicativo, il punteggio
resta \emph{trasparente e motivabile}: la spiegazione riporta i conteggi
grezzi per dimensione accanto al punteggio.

\begin{table}[ht]
\centering
\caption{Preset di priorit\`a e pesi associati.}
\label{tab:weights}
\small
\begin{tabular}{@{}lrrrr@{}}
\toprule
\textbf{Preset} & \textbf{Tropi} & \textbf{Regia} & \textbf{Cast} & \textbf{Generi} \\
\midrule
bilanciata & 3 & 2 & 2 & 5 \\
priorit\`a cast & 3 & 2 & 6 & 3 \\
priorit\`a regia & 3 & 6 & 2 & 3 \\
priorit\`a tropi & 6 & 2 & 2 & 3 \\
priorit\`a generi & 3 & 2 & 2 & 8 \\
\bottomrule
\end{tabular}
\end{table}

\subsection{Il grafo di spiegazione con CONSTRUCT}

La query \code{04\_explanation\_construct.rq} \`e il cuore della
spiegabilit\`a: produce un grafo RDF che \emph{racconta} la
raccomandazione. Il template \code{CONSTRUCT} emette la tripla
\code{:recommends}, la propriet\`a di condivisione opportuna
(\code{:sharesTropeWith}, \code{:sharesDirectorWith},
\code{:sharesActorWith}, \code{:sharesGenreWith}), i collegamenti alle
entit\`a condivise con le rispettive etichette e la provenienza tramite
\code{:assertedIn ?source}; il pattern \code{WHERE} \`e un \code{UNION}
di quattro rami, uno per dimensione, ciascuno dei quali interroga il
named graph della fonte:

\begin{lstlisting}[style=sparql, caption={Un ramo della query di
spiegazione (tropi condivisi con provenienza).}]
CONSTRUCT {
  <{{SEED}}> :recommends <{{REC}}> .
  <{{SEED}}> ?sharePredicate <{{REC}}> .
  <{{REC}}> :recommendedBecause ?trope .
  <{{SEED}}> :hasTrope ?trope .
  <{{REC}}> :hasTrope ?trope .
  <{{REC}}> :assertedIn ?source .
}
WHERE {
  {
    GRAPH ?source { <{{REC}}> :hasTrope ?trope }
    BIND (:sharesTropeWith AS ?sharePredicate)
    ?trope rdfs:label ?tropeLabel .
    ?source rdfs:label ?sourceLabel .
    FILTER (?source IN (:Wikidata, :DBpedia, :TMDB,
                        :TvTropes, :LocalAnalysis))
  }
  UNION { ... ramo regista ... }
  UNION { ... ramo cast ... }
  UNION { ... ramo generi ... }
}
\end{lstlisting}

Il risultato \`e un piccolo grafo (film seme $\rightarrow$ raccomandazione
$\rightarrow$ tropi/registi/attori/generi condivisi $\rightarrow$ fonti)
che il backend serializza in JSON e l'interfaccia disegna con vis.js.
Ogni \emph{perch\'e} della raccomandazione \`e un arco ispezionabile,
con la fonte a un clic. Si noti l'uso di \code{BIND} per selezionare la
propriet\`a di condivisione in funzione del ramo \code{UNION}: una
tecnica che tiene la query dichiarativa pura.

\subsection{Le altre query e i pattern ricorrenti}

Le restanti query coprono i pattern del corso:

\begin{itemize}
  \item \code{02\_shared\_tropes.rq} e \code{10\_cq\_shared\_two\_tropes.rq}:
        classificazione per sovrapposizione con \code{COUNT(DISTINCT)},
        \code{GROUP BY}, \code{ORDER BY DESC} e \code{HAVING (COUNT >= 2)}
        --- esattamente la forma suggerita nel testo di progetto;
  \item \code{11\_cq\_director\_trope.rq}: percorso semantico attraverso due
        relazioni (\code{directedBy} $\wedge$ \code{hasTrope}), verifica
        cq2;
  \item \code{12\_cq\_co\_stars.rq}: sfrutta direttamente
        \code{coStarsWith}; essendo la relazione completa (esplicita +
        inferita) nel grafo, la query non deve espandere la simmetria;
  \item \code{09\_genre\_hierarchy.rq}: gerarchia dei generi via
        \code{subGenreOf};
  \item \code{08\_source\_stats.rq}: statistiche per named graph con
        \code{GRAPH ?g}, per rendere verificabile il contributo di ogni
        fonte (cq8);
  \item \code{06\_mindbender.rq} e \code{07\_frequent\_collaborators.rq}:
        interroga classi e propriet\`a che \emph{esistono solo} nel grafo
        \code{:Inferred} --- uso effettivo dell'inferenza, non dei soli
        dati espliciti.
\end{itemize}

\subsection{Federazione SERVICE in ETL}

La federazione \`e confinata nell'ETL, come raccomandato dal testo di
progetto. La query \code{enrich\_birthdates\_service.rq} interroga
l'endpoint DBpedia per le date di nascita dei registi:

\begin{lstlisting}[style=sparql, caption={Arricchimento federato:
estrazione delle date di nascita da DBpedia.}]
PREFIX dbo: <http://dbpedia.org/ontology/>
PREFIX dbr: <http://dbpedia.org/resource/>
PREFIX : <http://example.org/cinetrope#>

SELECT ?person ?birthDate WHERE {
  ?person a :Regista .
  SERVICE <https://dbpedia.org/sparql> {
    ?dbpediaPerson dbo:birthDate ?birthDate .
    ?person owl:sameAs ?dbpediaPerson .
  }
}
\end{lstlisting}

Il risultato \`e salvato in \code{data/raw/dbpedia\_birthdates.json} e
ricaricato nel named graph \code{:DBpedia} al build successivo: la demo
usa solo il dato locale, rendendo il sistema riproducibile e immune
dall'indisponibilit\`a degli endpoint pubblici.

% ============================================================
\section{Applicazione web, API e CLI}\label{sec:web}
% ============================================================

\subsection{Stack applicativo}

Il back-end \`e un'applicazione FastAPI con template Jinja2; il front-end
usa vis.js per il grafo interattivo della spiegazione. Al primo avvio la
web app costruisce automaticamente lo store se
\code{data/generated/store.trig} non esiste (ETL + reasoner + validazione
+ scrittura del report). Le rotte principali sono in
tabella~\ref{tab:rotte}.

\begin{table}[ht]
\centering
\caption{Rotte della web app.}
\label{tab:rotte}
\footnotesize
\begin{tabular}{@{}lp{10.2cm}@{}}
\toprule
\textbf{Rotta} & \textbf{Descrizione} \\
\midrule
\code{/} & catalogo, statistiche del grafo, selezione del film \\
\code{/recommendations/\{film\}} & cinque raccomandazioni motivate con
preset di priorit\`a (\code{?priority=cast|director|tropes|genres}) \\
\code{/explanation/\{seed\}/\{rec\}} & grafo interattivo della spiegazione
(vis.js) con triple e serializzazione Turtle \\
\code{/queries} & competency questions con risultati e testo della query \\
\code{/sources} & statistiche e ruolo dei named graph \\
\code{/health} & stato sintetico del servizio \\
\code{/api/films} & catalogo in JSON \\
\code{/api/recommendations/\{seed\}} & raccomandazioni in JSON \\
\code{/api/explanation/\{seed\}/\{rec\}} & grafo di spiegazione in JSON \\
\code{/api/provenance/\{film\}/\{trope\}} & provenienza di un'asserzione \\
\code{/api/competency-questions} & tutte le CQ in JSON \\
\code{/api/stats} & statistiche del knowledge graph \\
\code{/docs} & documentazione OpenAPI generata \\
\bottomrule
\end{tabular}
\end{table}

\subsection{CLI e test}

La CLI (\code{python -m cinetrope ...}) espone \code{list},
\code{recommend}, \code{explain}, \code{cq}, \code{stats} e l'opzione
\code{--rebuild}, utile per rigenerare lo store dal terminale. La suite
\code{pytest} copre: costruzione del catalogo e dei named graph,
regole SWRL, raccomandazioni (inclusi i preset), grafo di spiegazione,
provenienza, controlli di validazione, competency questions,
arricchimento DBpedia e rotte web --- il che rende la pipeline verificabile
a ogni modifica dell'ontologia o delle query.

% ============================================================
\section{Valutazione e limiti}\label{sec:valutazione}
% ============================================================

\subsection{Risultati raggiunti}

Il sistema soddisfa tutti gli obiettivi della Sezione~1: il dataset
materializzato conta 500 film, 1 serie TV, 3.929 persone, 3.514
tropi/keyword e 21 generi, per 161.843 triple in otto grafi; il
ragionamento produce 52.608 inferenze; i cinque controlli di validazione
sono tutti a zero violazioni; le nove CQ sono eseguibili e la demo web
funziona offline. La tabella~\ref{tab:risultati} riassume gli indicatori
del report di build.

\begin{table}[ht]
\centering
\caption{Indicatori sintetici del build corrente.}
\label{tab:risultati}
\small
\begin{tabular}{@{}lr@{}}
\toprule
\textbf{Indicatore} & \textbf{Valore} \\
\midrule
Film / Serie TV & 500 / 1 \\
Persone / Tropi / Generi & 3.929 / 3.514 / 21 \\
Triple totali (named graph) & 161.843 \\
Triple inferite (\code{:Inferred}) & 52.608 \\
Inferenze SWRL (\code{hasFrequentCollaborator}) & 92 \\
Coppie di co-star inferite (simmetria) & 20.807 \\
\code{MindBender} inferiti & 0 \\
Violazioni di validazione & 0 \\
\bottomrule
\end{tabular}
\end{table}

\subsection{Limiti e scelte consapevoli}

La valutazione onesta richiede di dichiarare i limiti, alcuni dei quali
sono \emph{scelte} consapevoli, altri \emph{vincoli} da lavoro futuro:

\begin{itemize}
  \item \textbf{tropi = keyword TMDB}: tutte le keyword importate sono
        trattate come tropi condivisibili per massimizzare la copertura del
        recommender. Questo aumenta il richiamo ma abbassa la precisione
        semantica: \emph{dreams} o \emph{revenge} sono descrittori pi\`u
        deboli dei tropi narrativi canonici di TV Tropes. La classe
        \code{ElementoDescrittivo} \`e pronta per una curation futura;
  \item \textbf{MindBenderRule non innesca nel dataset corrente}: la regola
        fa riferimento agli individui \code{:UnreliableNarrator} e
        \code{:TwistEnding} del dataset curato originario; l'importazione
        TMDB usa identificatori diversi (slug minuscoli come
        \code{twistending}, e nessuna keyword \code{unreliablenarrator}
        \`e presente), quindi la regola produce 0 inferenze sul catalogo
        attuale. La regola resta corretta e testata; il mismatch \`e un
        esempio concreto del problema di \emph{conciliazione dei termini}
        visto a lezione: la soluzione naturale \`e una regola di
        allineamento (per esempio una \code{owl:equivalentProperty} o un
        mapping ETL fra slug e vocabolario curato);
  \item \textbf{motore SWRL parziale}: supporta solo \code{ClassAtom} e
        \code{IndividualPropertyAtom}; niente aggregazioni (per progettazione
        delegate all'ETL) e niente atomi built-in/SQWRL;
  \item \textbf{OWL RL e cardinalit\`a}: il profilo non implementa tutte
        le cardinalit\`a di OWL 2 DL; il progetto compensa con controlli
        SPARQL espliciti, accettando il costo di mantenerli separati dagli
        assiomi;
  \item \textbf{euristica trasparente, non appresa}: il punteggio \`e una
        combinazione lineare di conteggi con pesi dichiarati; non c'\`e
        apprendimento, e la qualit\`a dipende dalla copertura di keyword,
        cast e generi su TMDB;
  \item \textbf{scala}: la demo \`e ottimizzata per l'uso locale; un
        triple store con endpoint SPARQL (es.\ Fuseki) e reasoner
        distribuito sarebbero necessari per dataset di dimensione
        superiore.
\end{itemize}

\subsection{Lavoro futuro}

Le direzioni naturali di estensione sono: l'allineamento del vocabolario
dei tropi TMDB con TV Tropes tramite \code{owl:sameAs} e
\code{owl:equivalentProperty}; la migrazione da rdflib a un triple store
con endpoint SPARQL per query federate \emph{reali} a runtime; il supporto
SQWRL per aggregazioni dichiarative; l'uso della classe
\code{Scena}/\code{annotatedWith} per annotazioni a livello di scena; e la
validazione dell'ontologia con reasoner OWL DL completi (HermiT/Pellet)
oltre al profilo RL.

% ============================================================
\section{Conclusioni}
% ============================================================

CineTrope dimostra che la teoria del Semantic Web --- presentata a lezione
come stack di linguaggi sempre pi\`u espressivi --- si traduce in un
sistema applicativo \emph{funzionante e spiegabile per costruzione}. Ogni
concetto teorico ha un corrispettivo osservabile: le triple RDF e i named
graph sono il formato del dataset TriG; RDFS e OWL definiscono la
tassonomia di classi e le propriet\`a inverse, simmetriche ed equivalenti
che il reasoner \code{owlrl} espande in 52.608 triple inferite; SWRL ---
con il compromesso aggregazioni/ETL discusso a lezione --- produce i
collegamenti di collaborazione frequente; SPARQL \code{SELECT} calcola le
raccomandazioni, \code{CONSTRUCT} genera il grafo della spiegazione e
\code{GRAPH ?source} rende verificabile la provenienza di ogni fatto; la
query federata \code{SERVICE} arricchisce il grafo in fase ETL senza
comprometterne la riproducibilit\`a.

La lezione metodologica pi\`u importante \`e forse la complementarit\`a dei
livelli: nessun singolo strumento basta. Le cardinalit\`a chiedono OWL, ma
il profilo RL le copre solo in parte e servono controlli espliciti; le
regole SWRL collegano individui, ma non contano; l'aggregazione sta in
SPARQL; la provenienza sta nei named graph, non nella reificazione; e
l'Open World Assumption ricorda che i controlli \emph{segnalano} senza
\emph{correggere}. Saper orchestrare questi compromessi --- pi\`u che
conoscere i singoli costrutti --- \`e ci\`o che il progetto ha insegnato.

% ============================================================
\begin{thebibliography}{9}
% ============================================================

\bibitem{allemang2020}
D.~Allemang, J.~Hendler, \emph{Semantic Web for the Working Ontologist},
3$^{\mathrm{a}}$ ed., Elsevier, 2020.

\bibitem{antoniou2012}
G.~Antoniou, P.~Groth, F.~van Harmelen, R.~Hoekstra, \emph{A Semantic Web
Primer}, 3$^{\mathrm{a}}$ ed., MIT Press, 2012.

\bibitem{rdf11}
W3C, \emph{RDF 1.1 Concepts and Abstract Syntax}, Recommendation, 2014.
\url{https://www.w3.org/TR/rdf11-concepts/}

\bibitem{rdfs}
W3C, \emph{RDF Schema 1.1}, W3C Recommendation, 2014.
\url{https://www.w3.org/TR/rdf-schema/}

\bibitem{owl2}
W3C OWL Working Group, \emph{OWL 2 Web Ontology Language Document
Overview}, 2$^{\mathrm{a}}$ ed., 2012.
\url{https://www.w3.org/TR/owl2-overview/}

\bibitem{owl2profiles}
B.~Motik, B.~Cuenca Grau, I.~Horrocks, et al., \emph{OWL 2 Web Ontology
Language Profiles}, W3C Recommendation, 2012.
\url{https://www.w3.org/TR/owl2-profiles/}

\bibitem{sparql11}
W3C, \emph{SPARQL 1.1 Query Language}, Recommendation, 2013.
\url{https://www.w3.org/TR/sparql11-query/}

\bibitem{swrl}
I.~Horrocks, P.~F.~Patel-Schneider, H.~Boley, et al., \emph{SWRL: A
Semantic Web Rule Language Combining OWL and RuleML}, W3C Member
Submission, 2004. \url{https://www.w3.org/Submission/SWRL/}

\bibitem{studer98}
R.~Studer, V.~R.~Benjamins, D.~Fensel, \emph{Knowledge Engineering:
Principles and Methods}, in \emph{Data and Knowledge Engineering},
25(1--2), 1998.

\bibitem{rdflib}
rdflib contributors, \emph{rdflib --- a Python library for working with
RDF}. \url{https://github.com/RDFLib/rdflib}

\bibitem{owlrl}
N.~Bassiliades, et al./owl-rl contributors, \emph{OWL-RL --- a simple OWL
RL reasoner in Python}. \url{https://github.com/RDFLib/owlrl}

\end{thebibliography}

\end{document}
